\chapter*{Liste des acronymes}
\addcontentsline{toc}{chapter}{Liste des acronymes}

\section*{Institutions, organismes, associations et projets}

\begin{itemize}
	\item \textbf{Abes} : Agence bibliographique de l'enseignement supérieur
	\item \textbf{ABF} : Association des bibliothécaires de France
	\item \textbf{ACIM} : Association pour la coopération des professionnels de l'information musicale
	\item \textbf{ADBU} : Association des directeurs et personnels de direction des bibliothèques universitaires et de la documentation
	\item \textbf{ANR} : Agence nationale de la recherche
	\item \textbf{BBF} : \textit{Bulletin des bibliothèques de France}
	\item \textbf{BnF} : Bibliothèque nationale de France
	\item \textbf{BPI} : Bibliothèque publique d'information
	\item \textbf{BU} : Bibliothèque universitaire
	\item \textbf{CFA} : Centre de formation d'apprentis
	\item \textbf{CMBV} : Centre de musique baroque de Versailles
	\item \textbf{CMF} : Confédération musicale de France
	\item \textbf{CNFPT} : Centre national de la fonction publique territoriale
	\item \textbf{CNSMD} : Conservatoire national supérieur de musique et de danse (Paris et Lyon)
	\item \textbf{CNSMDP} : Conservatoire national supérieur de musique et de danse de Paris
	\item \textbf{CRFCB} : Centres régionaux de formation aux métiers des bibliothèques
	\item \textbf{CSB} : Comité stratégique bibliographique
	\item \textbf{Doremus} : \textit{Doing Reusable Music Data} (projet ANR)
	\item \textbf{ENSSIB} : École nationale supérieure des sciences de l'information et des bibliothèques
	\item \textbf{HEM} : Haute École de musique (Genève)
	\item \textbf{IBM} : \textit{International Business Machines Corporation}
	\item \textbf{IFLA} : Fédération internationale des associations et institutions de bibliothèques (\textit{International Federation of Library Associations and Institutions})
	\item \textbf{IReMus} : Institut de recherche en musicologie (Unité mixte de recherche)
	\item \textbf{MMP} : Médiathèque Musicale de Paris
	\item \textbf{PEPR-Iccare} : Programme d'équipements prioritaires de recherche - Industries culturelles et créatives :  action, recherche, expérimentation (projet ciblé ANR)
	\item \textbf{PUR} : Presses universitaires de Rennes 
	\item \textbf{UNESCO} : Organisation des Nations unies pour l'éducation, la science et la culture (\textit{United Nations Educational, Scientific and Cultural Organization})
	\item \textbf{W3C} : \textit{World Wide Web Consortium}
\end{itemize}

\section*{Catalogues, bases de données, identifiants et Web de données}

\begin{itemize}
	\item \textbf{ARK} : \textit{Archival Resource Key}
	\item \textbf{CCFr} : Catalogue collectif de France
	\item \textbf{HAL} : Hyper articles en ligne (plateforme d'archives ouvertes)
	\item \textbf{IdRef} : Identifiants et référentiels pour l'enseignement supérieur et la recherche
	\item \textbf{ISNI} : \textit{International Standard Name Identifier}
	\item \textbf{MIMO} : \textit{Musical Instrument Museums Online} (base de données internationale pour les instruments de musique)
	\item \textbf{OWL} : \textit{Web Ontology Language}
	\item \textbf{RDF} : \textit{Resource Description Framework}
	\item \textbf{RISM} : Répertoire international des sources musicales (\textit{International Inventory of Musical Sources})
	\item \textbf{SKOS} : \textit{Simple Knowledge Organization System}
	\item \textbf{Sudoc} : Système universitaire de documentation
	\item \textbf{URI} : \textit{Uniform Resource Identifier}
	\item \textbf{URL} : \textit{Uniform Resource Locator}
	\item \textbf{VIAF} : \textit{Virtual International Authority File}
\end{itemize}

\section*{Modèles de données, normes bibliothéconomiques et classifications}

\begin{itemize}
	\item \textbf{4 C} : Collecter, classer, conserver, communiquer (principe d'archivistique)
	\item \textbf{CDU} : Classification décimale universelle
	\item \textbf{FRBR} : \textit{Functional Requirements for Bibliographic Records}
	\item \textbf{IFLA-LRM / LRM} : \textit{Library Reference Model}
	\item \textbf{ISO} : \textit{International Organization for Standardization}
	\item \textbf{MARC} : \textit{Machine-Readable Cataloging}
	\item \textbf{MARC21} : \textit{MARC 21st Century}
	\item \textbf{OEMI} : Œuvre, Expression, Manifestation, Item (déclinaison française du modèle LRM / FRBR)
	\item \textbf{PCDM / PCDM4} : Principe de classement des documents musicaux (4\textsuperscript{e} version)
	\item \textbf{RAMEAU} : Répertoire d'autorité matière encyclopédique et alphabétique unifié
	\item \textbf{RDA-FR} : \textit{Resource Description and Access} — version française
	\item \textbf{SPAR} : Système de préservation et d'archivage réparti
	\item \textbf{UNIMARC} : \textit{Universal MARC}
	\item \textbf{UNIMARC ER} : UNIMARC Entités-Relations
\end{itemize}

\section*{Informatique, formats, outils et systèmes d'information}

\begin{itemize}
	\item \textbf{CSV} : \textit{Comma-Separated Values}
	\item \textbf{IA} : Intelligence artificielle
	\item \textbf{JSON} : \textit{JavaScript Object Notation}
	\item \textbf{MusicXML} : \textit{Music Extensible Markup Language} (format d'échange pour la notation musicale)
	\item \textbf{PDF} : \textit{Portable Document Format}
	\item \textbf{PMB} : Logiciel libre de gestion de bibliothèque
	\item \textbf{SGBDR} : Système de gestion de bases de données relationnelles
	\item \textbf{SI} : Système d'information
	\item \textbf{SIGB} : Système intégré de gestion de bibliothèque
	\item \textbf{SQL} : \textit{Structured Query Language}
	\item \textbf{UX} : \textit{User Experience} (expérience utilisateur)
	\item \textbf{XML} : \textit{Extensible Markup Language}
\end{itemize}

\section*{Métiers, diplômes, concepts et pratiques}

\begin{itemize}
	\item \textbf{DCB} : Diplôme de conservateur des bibliothèques
	\item \textbf{EMI} : Éducation aux médias et à l'information
	\item \textbf{HIP} : \textit{Historically Informed Performance} / \textit{Historically Informed Practice} (pratique / interprétation historiquement informée)
	\item \textbf{IOUPI} : Incorrect, Ordinaire, Périmé, Usé, Inadéquat (méthode de désherbage)
	\item \textbf{IST} : Information scientifique et technique
	\item \textbf{LMD} : Licence, Master, Doctorat (schéma européen de l'enseignement supérieur)
	\item \textbf{PEB} : Prêt entre bibliothèques
	\item \textbf{Poldoc} : Politique documentaire 
\end{itemize}

\section*{Musique, supports et abbréviations matérielles}

\begin{itemize}
	\item \textbf{CD} : \textit{Compact Disc}
	\item \textbf{MAO} : Musique assistée par ordinateur
	\item \textbf{Pno.} : Piano (abréviation courante sur les parties séparées et le matériel d'orchestre)
\end{itemize}